% !TeX root = ../教材题库_全部条目_转写工作稿.tex
% Source: XJ-B1-C04-TH-P098-Y0419
% 第4章 / 4.1 实数指数幂和幂函数 / 教材 p.98 / PDF 第105页
% OCR未校对：公式与文字均可能有误；未完成数学TeX转写；禁止直接发给学生
% 原题及相邻原解截图（核对用）：../题目原图/XJ-B1-C04-TH-P098-Y0419-01.webp
% 本文件只在转写完成后作为题干使用；OCR 条目尚不是可靠题干。
\begin{bankitem}{p.98 思考（含紧邻说明）}{第4章 / 4.1 实数指数幂和幂函数 / 教材 p.98 / PDF 第105页}{XJ-B1-C04-TH-P098-Y0419}{OCR未校对：公式与文字均可能有误；未完成数学TeX转写；禁止直接发给学生}
\begin{ocrdraft}
[教材p.98 / PDF第105页]
\par
？思考我们知道，对任意的正整数刀和正数若磊>1，则>1；若磊<1，
\par
则a"<l.那么，对于任意的正有理数，和1之间是否也有类似的关系呢？
\par
设a>1，可以猜测下面用反证法证明这一猜测成立．
\par
假设<1，
\par
则（层)"=a<l,与已知>1矛盾，故产>1.
\par
当正数a<I时，<
\par
进一步可推知，
\par
I成立吗？
\par
综合起来得到有关有理数指数幂的基本不等式：
\par
对任意的正有理数厂和正数a，
\par
若磊>1，则>1；若磊<1，则d<l.
\par
根据负指数的意义和倒数的性质可推得：
\par
对任意的负有理数厂和正数a，
\par
若磊>1，则<1；若磊<1，则d>l.
\par
由此可知：
\par
对任意的正数磊>1和两有理数厂>s，有a一磊一>1，
\par
即>．
\par
对任意的正数磊<1和两有理数厂>s，有a
\par
即<．
\par
这两个不等式，有助于建立和理解无理数指数幂的概念
\par
0
\par
无理数指数幂的概念
\end{ocrdraft}
\end{bankitem}
